\documentclass[11pt,a4paper]{article}

\usepackage[T1]{fontenc}
\usepackage[utf8]{inputenc}
\usepackage{lmodern}
\usepackage{amssymb,amsmath,amsthm,mathtools}
\usepackage[margin=2.5cm]{geometry}
\usepackage{parskip}
\usepackage[hidelinks]{hyperref}
\usepackage{microtype}

% Theorem environments (arXiv / journal style)
\theoremstyle{plain}
\newtheorem{theorem}{Theorem}[section]
\newtheorem{proposition}[theorem]{Proposition}
\newtheorem{lemma}[theorem]{Lemma}
\newtheorem{corollary}[theorem]{Corollary}

\theoremstyle{definition}
\newtheorem{definition}[theorem]{Definition}
\newtheorem{example}[theorem]{Example}
\newtheorem{remark}[theorem]{Remark}

% Notation
\newcommand{\Shape}{\mathrm{Shape}}
\newcommand{\Pos}{\mathrm{Pos}}
\newcommand{\Type}{\mathbf{Type}}
\newcommand{\id}{\mathrm{id}}
\newcommand{\src}{\mathrm{src}}
\newcommand{\tgt}{\mathrm{tgt}}
\newcommand{\rootpos}{\mathrm{root}}
\newcommand{\sub}{\mathrm{sub}}
\newcommand{\shift}{\mathrm{shift}}

\title{Polytrees and Thinness of Network Directed Containers}
\author{Extracted from a machine-checked development}
\date{}

\begin{document}
\maketitle

\begin{abstract}
A directed graph is a \emph{polytree} when it is acyclic and satisfies an
edge-disjoint reachability condition: distinct edges out of the same source
never share a common reachable descendant. We prove that if $G$ is a polytree,
then between any two nodes there is at most one directed path---equivalently,
the network directed container $\mathrm{network}(G)$ is thin. As a corollary,
positions at a node~$v$ are in bijection with the set of nodes reachable
from~$v$, the unique path serving as the inverse to the endpoint projection.
The proofs of all stated results have been machine-checked; we record only the
mathematical content.
\end{abstract}

\tableofcontents
\bigskip

%======================================================================
\section{Introduction}
%======================================================================

The network construction associates to every directed graph~$G$ a directed
container whose positions at a node are directed paths out of that node,
paired with their endpoints. In general there may be many paths between a
given pair of nodes. When $G$ is a polytree---the directed analogue of a
forest---paths become unique, and the position type collapses to a mere set of
reachable endpoints.

The present note isolates that uniqueness theorem. We take as definition of
polytree the conjunction of directed acyclicity with an edge-disjointness
condition on reachability: if two distinct edges leave the same vertex, no
vertex is reachable from both of their targets. From these two hypotheses one
obtains, by induction on paths, that any two paths with the same endpoints
coincide. Thinness of the network container and a bijection between positions
and reachable nodes are immediate consequences.

The development works throughout in a single universe of types. Proofs are
indicated by the phrase ``by lean,'' meaning they have been checked by a proof
assistant rather than written out in prose.

\paragraph{What this note claims.}
In ordinary mathematical language: polytrees have unique directed paths between
any ordered pair of nodes. Consequently the free category on a polytree is a
preorder (at most one morphism between any two objects), and the directed
container of routes is thin.

%======================================================================
\section{Reachability, acyclicity, and edge-disjointness}
%======================================================================

Throughout, let $G$ be a directed graph with node type~$V$, edge type~$E$, and
endpoint maps $\src,\tgt : E \to V$. Write $\mathrm{Path}_G(v,w)$ for the type
of finite directed paths from $v$ to~$w$, as in the network development.

\begin{definition}[Reachability]\label{def:reachable}
A node $w$ is \textbf{reachable} from a node~$v$ when the type
$\mathrm{Path}_G(v,w)$ is inhabited. We write $\mathrm{Reachable}_G(v,w)$ for
this proposition.
\end{definition}

\begin{definition}[Trivial path]\label{def:trivial}
A path $p : \mathrm{Path}_G(v,w)$ is \textbf{trivial} when it is the empty path
(equivalently, when $v = w$ and $p = \mathrm{nil}(v)$). A non-empty path is a
$\mathrm{cons}$.
\end{definition}

\begin{definition}[Acyclicity]\label{def:acyclic}
The graph $G$ is \textbf{acyclic} when every self-path is trivial:
\begin{equation}
  \forall\, v \in V.\;
  \forall\, p : \mathrm{Path}_G(v,v).\;
  \text{$p$ is trivial}.
\end{equation}
\end{definition}

\begin{definition}[Edge-disjoint reachability]\label{def:edge-disjoint}
The graph $G$ satisfies \textbf{edge-disjoint reachability} when, for any two
edges $e_1 \neq e_2$ with $\src(e_1) = \src(e_2)$ and any node~$w$,
\begin{equation}
  \neg\Bigl(
    \mathrm{Reachable}_G\bigl(\tgt(e_1),\, w\bigr)
    \;\wedge\;
    \mathrm{Reachable}_G\bigl(\tgt(e_2),\, w\bigr)
  \Bigr).
\end{equation}
In words: distinct edges out of the same source never have a common reachable
descendant.
\end{definition}

\begin{definition}[Polytree]\label{def:polytree}
A \textbf{polytree} is a directed graph that is acyclic and satisfies
edge-disjoint reachability.
\end{definition}

\begin{lemma}[No reverse cons-path]\label{lem:no-cons}
If $G$ is acyclic, $e \in E$, and $p : \mathrm{Path}_G\bigl(\tgt(e),\,
\src(e)\bigr)$, then contradiction: $\mathrm{cons}(e,p)$ would be a non-trivial
self-path at $\src(e)$.
\end{lemma}

\begin{proof}
By lean.
\end{proof}

%======================================================================
\section{Path uniqueness}
%======================================================================

To compare two paths without nested dependent elimination on indices, it is
convenient to package the first-step decomposition of a path explicitly.

\begin{definition}[Path view]\label{def:path-view}
For nodes $v,w \in V$, a \textbf{path view} is either
\begin{enumerate}
  \item $\mathrm{trivial}(h)$ for a proof $h : v = w$, or
  \item $\mathrm{step}(e,h_{\mathrm{src}},\tau)$ for an edge~$e$, a proof
        $h_{\mathrm{src}} : \src(e) = v$, and a tail path
        $\tau : \mathrm{Path}_G\bigl(\tgt(e),\, w\bigr)$.
\end{enumerate}
Every path $p : \mathrm{Path}_G(v,w)$ determines a path view, and every path
view determines a path; the two constructions are mutually inverse on paths.
\end{definition}

\begin{theorem}[Path uniqueness in a polytree]\label{thm:path-unique}
Let $G$ be a polytree. For any nodes $v,w \in V$ and any two paths
$p,q : \mathrm{Path}_G(v,w)$, one has $p = q$.
\end{theorem}

\begin{proof}
By lean.
\end{proof}

\begin{remark}[Idea of the argument]
The proof proceeds by structural induction on~$p$, examining the path view
of~$q$ at each step.
\begin{itemize}
  \item If $p$ is empty and $q$'s view is trivial, both paths are empty and
        agree.
  \item If $p$ is empty and $q$'s view is a step, the step produces a directed
        cycle, contradicting acyclicity (Lemma~\ref{lem:no-cons}).
  \item If $p$ is a cons and $q$'s view is trivial, one obtains a cycle
        symmetrically.
  \item If $p = \mathrm{cons}(e_1,p_1)$ and $q$'s view is a step along~$e_2$:
        when $e_1 \neq e_2$, the common endpoint~$w$ is reachable from both
        $\tgt(e_1)$ and $\tgt(e_2)$, contradicting edge-disjointness; when
        $e_1 = e_2$, the inductive hypothesis equates the tails.
\end{itemize}
\end{remark}

\begin{corollary}[Thinness]\label{cor:thin}
If $G$ is a polytree, then for every pair of nodes $v,w$ the type
$\mathrm{Path}_G(v,w)$ is a subsingleton: it has at most one element.
Equivalently, the directed container $\mathrm{network}(G)$ is thin.
\end{corollary}

\begin{proof}
By lean.
\end{proof}

%======================================================================
\section{Positions and reachable endpoints}
%======================================================================

In a polytree, a position at~$v$---a pair $\langle w,\rho \rangle$ of endpoint
and path---carries no more information than the endpoint alone, once
reachability is known.

\begin{definition}[Endpoint equivalence]\label{def:endpoint-equiv}
Let $G$ be a polytree and fix a node~$v$. Define maps
\begin{align}
  \Phi
  &\;:\;
  \sum_{w \,:\, V}\; \mathrm{Path}_G(v,w)
  \;\longrightarrow\;
  \bigl\{ w \in V \mid \mathrm{Reachable}_G(v,w) \bigr\}, \\[4pt]
  \Phi\bigl\langle w,\rho \bigr\rangle
  &\;=\;
  \bigl\langle w,\, \langle\rho\rangle \bigr\rangle,
\end{align}
and, choosing for each reachability witness a path (via the axiom of choice),
\begin{align}
  \Psi
  &\;:\;
  \bigl\{ w \in V \mid \mathrm{Reachable}_G(v,w) \bigr\}
  \;\longrightarrow\;
  \sum_{w \,:\, V}\; \mathrm{Path}_G(v,w), \\[4pt]
  \Psi\bigl\langle w,\, h \bigr\rangle
  &\;=\;
  \bigl\langle w,\, \mathrm{choose}(h) \bigr\rangle.
\end{align}
\end{definition}

\begin{theorem}[Endpoint bijection]\label{thm:endpoint-equiv}
If $G$ is a polytree, then for every node~$v$ the maps $\Phi$ and $\Psi$ of
Definition~\ref{def:endpoint-equiv} are mutually inverse. In particular,
\begin{equation}
  \Pos\bigl(\mathrm{network}(G)\bigr)(v)
  \;\simeq\;
  \bigl\{ w \in V \mid \mathrm{Reachable}_G(v,w) \bigr\}.
\end{equation}
\end{theorem}

\begin{proof}
By lean.
\end{proof}

\begin{remark}
The left inverse uses path uniqueness (Theorem~\ref{thm:path-unique}) to
identify the chosen path with the original path component. The right inverse is
immediate from the construction of $\Phi$ and $\Psi$. Choice enters only in
selecting a path from a mere inhabitance proof; uniqueness ensures the choice
is immaterial.
\end{remark}

%======================================================================
\section{Discussion}
%======================================================================

\paragraph{What has been established.}
Polytrees---acyclic digraphs with edge-disjoint reachability---have unique
directed paths. The network directed container is therefore thin, and positions
at a node biject with reachable endpoints. In supply-chain terms: on a
polytree shipping graph, a destination determines a unique route from a given
origin.

\paragraph{On the definition of polytree.}
The edge-disjointness condition is the directed counterpart of the undirected
forest property that two distinct branches never reconverge. Acyclicity alone
does not suffice for uniqueness (a directed acyclic graph may have multiple
paths between the same endpoints); both hypotheses are used in the inductive
argument.

\paragraph{Scope.}
The development does not construct the polytree structure on the running
eleven-facility wine network (though that graph is a polytree by inspection),
nor does it treat the converse: thinness of $\mathrm{network}(G)$ characterising
polytrees. Counterexamples on the DAG and cyclic rungs of the ladder appear
separately.

\paragraph{Machine checking.}
Every definition and theorem statement above corresponds to a declaration in a
machine-checked formalisation. The phrase ``by lean'' indicates that the proof
has been verified by a proof assistant.

%======================================================================
\section*{Acknowledgements}
%======================================================================

The mathematical content follows the thinness proposition for polytree networks
in the supply-chain development that motivated this note. The presentation is a
prose rendering of a formal development intended for human review of the
theorem statements.

\begin{thebibliography}{9}

\bibitem{abbott2005}
M.~Abbott, T.~Altenkirch, and N.~Ghani.
\newblock Containers: constructing strictly positive types.
\newblock \emph{Theoretical Computer Science}, 342(1):3--27, 2005.

\bibitem{ahman2014}
D.~Ahman and T.~Uustalu.
\newblock Coalgebraic update lenses.
\newblock In \emph{MFPS}, volume~308 of \emph{ENTCS}, pages 25--48, 2014.

\end{thebibliography}

\end{document}
